\documentclass[10pt,a4paper]{amsart}
\usepackage[margin=19mm]{geometry}
\usepackage{amsmath,amssymb,mathtools}
\usepackage[scr=rsfs,cal=euler]{mathalfa}
\usepackage{xfrac}
\usepackage[unicode,hidelinks,bookmarks=false]{hyperref}
\newcommand{\ie}{i.e.\ }
\newcommand{\cf}{cf.\ }
\newcommand{\resp}{respectively\ }
\newcommand{\Subsection}{\subsection}
\providecommand{\nobreakdash}{\nobreak\hbox{-}}
\setlength{\emergencystretch}{2em}
\multlinegap0pt
\usepackage{amssymb}
\usepackage{mathtools}
\usepackage[shortlabels]{enumitem}
\usepackage{tikz}
\newtheorem{proposition}{Proposition}
\newcommand{\Hh}{\mathcal H}
\newcommand{\I}{\mathbf I}
\newcommand{\M}{\mathbf M}
\newcommand{\llbracket}{[\![}
\newcommand{\rrbracket}{]\!]}
\title{\bf Odd branching obstructs multiplicity-one integral current structures}
\author{Deguang Zhong}
\date{Source: arXiv:2609.14515v1 (CC BY 4.0)}
\begin{document}
\maketitle

\noindent\emph{Source and licence.} \bf Odd branching obstructs multiplicity-one integral current structures. Version arXiv:2609.14515v1 (CC BY 4.0), distributed by arXiv under the Creative Commons Attribution 4.0 International licence, http://creativecommons.org/licenses/by/4.0/, as recorded in the arXiv licence metadata for this exact version and shown on the abstract page https://arxiv.org/abs/2609.14515v1. Full licence notice: \texttt{licenses/arxiv-2609.14515-CC-BY-4.0.txt}.\par\smallskip

\noindent\textit{Editorial scope.} Counted unit: the article's proposition prop:mod2 (source0.tex lines 289--313). The article's complete original proof of it is delivered with it (source0.tex lines 315--381, one \texttt{proof} environment of 67 lines). No mathematics is written, repaired or re-derived: the statement and the proof are the article's own lines, in the article's own notation, and no retained line is substituted or re-flowed.  No further source-local context is needed: every cross-reference of every retained line resolves inside the two delivered spans.  Nothing was omitted from any retained span, so the package carries no omission record.  No retained line contains a citation, so the wrapper carries no bibliography and no bibliography entry.  No retained line contains a figure environment, an image-inclusion command or drawing code, so no image is delivered and the package declares no asset dependency.  Editorial formatting only, in addition: an AMS \texttt{amsart} preamble that restates the article's own declarations and macros used by the retained lines, namely the environments \texttt{proposition} on separate counters, together with the standard packages those lines need.  The retained environments print in this document as Proposition 1 in order of appearance, whereas the article prints the same items with the numbering of its own full document.  The complete native source of the article as distributed in the arXiv e-print for this exact version is bundled unchanged as \texttt{sources/210324/source0.tex}.
\begin{proposition}[The canonical mod-$2$ obstruction]\label{prop:mod2}
Let $H$ be a graph-like space with countably many edges
$E(H)=\{e_1,e_2,\ldots\}$ such that
\[
       \sum_{e\in E(H)}\ell(e)<\infty,
\]
and suppose that the complement of the open edge interiors is
$\Hh^1$-null. Give each edge an arbitrary orientation and let
\[
       C_H=\sum_{e\in E(H)}[\llbracket e\rrbracket]_2,
\]
where the series is understood in mod-$2$ mass, and hence also in the
mod-$2$ flat norm. Then $C_H$ is a rectifiable flat chain modulo $2$ of
finite mod-$2$ mass and is independent of the chosen orientations. At every
isolated vertex $v$ of finite degree, the local boundary coefficient of
$C_H$ is
\[
       \deg_H(v)\pmod2.
\]
Consequently, if $H$ has infinitely many isolated odd-degree vertices, then
\[
       \M_2(\partial C_H)=\infty;
\]
in particular, $C_H$ is not an integral flat chain modulo $2$. Moreover, any $S\in\I_1(H)$ having absolute multiplicity one almost everywhere on every open edge would satisfy $[S]_2=C_H$.
\end{proposition}

\begin{proof}
For
\[
       C_N=\sum_{j=1}^{N}[\llbracket e_j\rrbracket]_2
\]
and $M>N$, subadditivity of the relaxed mass gives
\[
 \mathcal F_2(C_M-C_N)
 \le \M_2(C_M-C_N)
 \le \sum_{j=N+1}^{M}\ell(e_j).
\]
Thus $(C_N)$ is Cauchy in mod-$2$ mass and in the mod-$2$ flat norm; denote
its flat limit by $C_H$. The countable edge representation, together with
the disjointness of the open edge interiors and the nullity of their
complement, shows directly that $C_H$ is rectifiable and that
\[
       \M_2(C_H)\le \sum_{e\in E(H)}\ell(e)<\infty.
\]
More precisely, lower semicontinuity of $\M_2$ under flat convergence gives
\[
       \M_2(C_H-C_N)
       \le \sum_{j>N}\ell(e_j)\longrightarrow0,
\]
so the displayed series does converge in mod-$2$ mass as asserted.
Reversing the orientation of an edge changes its integer coefficient from
$1$ to $-1$, which has no effect modulo $2$. Hence the limit is independent
of all orientation choices.

The estimate
\[
       \mathcal F_2(\partial A)\le \mathcal F_2(A)
\]
shows that $\partial C_N\to\partial C_H$ in the mod-$2$ flat norm. Let $v$
be an isolated vertex of degree $d$. All the finitely many edges incident to
$v$ occur in some $C_N$. Choose a star neighborhood $U_v$ containing no
other vertex and meeting only those incident edges. On $U_v$ the tail of the
series vanishes, so the local boundary of $C_H$ agrees with that of the
finite sum of its $d$ incident edge chains. Each incident oriented edge has
coefficient $1$ or $-1$ at $v$, and both reduce to $1$ modulo $2$. Hence the
local coefficient of $\partial C_H$ at $v$ is $d$ modulo $2$.

Now let $F$ be any finite set of isolated odd-degree vertices. Shrinking the
star neighborhoods, they may be chosen pairwise disjoint. In each such
neighborhood $\partial C_H$ has a nonzero $0$-dimensional coefficient and
therefore mod-$2$ mass at least one. Consequently
\[
       \M_2(\partial C_H)\ge \#F.
\]
If there are infinitely many isolated odd-degree vertices, finite sets $F$
of arbitrarily large cardinality give
$\M_2(\partial C_H)=\infty$. Thus $C_H$ has no finite-mass integral boundary
modulo $2$ and is not an integral flat chain modulo $2$.

Finally, suppose that $S\in\I_1(H)$ has absolute multiplicity one almost
everywhere on every open edge. In the rectifiable representation of $S$ on
an edge, its signed integer coefficient is therefore $1$ or $-1$ almost
everywhere; it need not be constant if $\partial S$ has atoms in the edge
interior, but its reduction modulo $2$ is nevertheless identically one.
Thus $[S]_2$ and $C_H$ agree on the union of the open edge interiors. Both
one-dimensional mass measures vanish on the complementary $\Hh^1$-null
set, so $[S]_2=C_H$. Since reduction commutes with the boundary,
\[
       \partial C_H=\partial[S]_2=[\partial S]_2,
\]
whose mod-$2$ mass is finite because $\partial S\in\I_0(H)$. This recovers
the contradiction when infinitely many isolated odd vertices are present.
\end{proof}

\end{document}
